\documentclass[10pt,a4paper]{amsart}
\usepackage[margin=19mm]{geometry}
\usepackage{amsmath,amssymb,mathtools}
\usepackage[scr=rsfs,cal=euler]{mathalfa}
\usepackage{xfrac}
\usepackage[unicode,hidelinks,bookmarks=false]{hyperref}
\newcommand{\ie}{i.e.\ }
\newcommand{\cf}{cf.\ }
\newcommand{\resp}{respectively\ }
\newcommand{\Subsection}{\subsection}
\providecommand{\nobreakdash}{\nobreak\hbox{-}}
\setlength{\emergencystretch}{2em}
\multlinegap0pt
\usepackage{bm}
\usepackage{bbm}
\usepackage{dsfont}
\usepackage{xspace}
\usepackage{cancel}
\usepackage{mathtools}
\usepackage{amsthm}
\makeatletter
\@ifundefined{theorem}{\newtheorem{theorem}{Theorem}}{}
\@ifundefined{claim}{\newtheorem{claim}[theorem]{Claim}}{}
\@ifundefined{lemma}{\newtheorem{lemma}[theorem]{Lemma}}{}
\makeatother
\makeatletter
\renewcommand{\geq}{\geqslant}
\renewcommand{\leq}{\leqslant}
\expandafter\ifx\csname repdefinition\endcsname\relax
\newenvironment{repdefinition}[1]
  {\renewcommand\therdef{\ref*{#1}}\rdef}
\fi
\expandafter\ifx\csname replemma\endcsname\relax
\newenvironment{replemma}[1]
  {\renewcommand\therlemma{\ref*{#1}}\rlemma}
\fi
\expandafter\ifx\csname reptheorem\endcsname\relax
\newenvironment{reptheorem}[1]
  {\renewcommand\therthm{\ref*{#1}}\rthm}
  {\endrthm}
\fi
\makeatother
\title[Maximum Percolation Time on the q-ary Hypercube]{Maximum Percolation Time on the q-ary Hypercube}
\author{Fengxing Zhu}
\date{Source: arXiv:2503.00990v2 (CC BY 4.0)}
\begin{document}
\maketitle

\noindent\emph{Source and licence.} Maximum Percolation Time on the q-ary Hypercube. Version arXiv:2503.00990v2 (CC BY 4.0), distributed under the Creative Commons Attribution 4.0 International licence, as recorded in the arXiv licence metadata for this exact version and shown on the abstract page https://arxiv.org/abs/2503.00990v2. Full rights notice: licenses/arxiv-2503.00990-CC-BY-4.0.txt.\par\smallskip

\noindent\textit{Editorial scope.} Counted unit: the Theorem whose source label is \texttt{None} in the article (source0.tex lines 641--646, source label \texttt{None}), delivered together with the article's own complete original proof of it (source0.tex lines 647--951, one \texttt{proof} environment of 305 lines). No mathematics is written, repaired or re-derived: statement and proof are the article's own text line for line. Required context retained uncounted: lemma:nested sequence (lines 334--336); lemma:ST3 (lines 423--450); lemma:ST6 (lines 540--563); lemma:summary (lines 616--624). Every definition, display and label of those lines is retained unchanged. Omitted from this package and not redistributed: 1--333, 337--422, 451--539, 564--615, 625--640, 952--978. Nothing inside a retained excerpt is omitted or altered: every retained body is delivered exactly as the article's lines are written, so no omission receipt of any kind accompanies this package. Citations: the retained excerpts carry no citation command at all, so no bibliography entry of the article is transcribed and no reference is redistributed. Reference adaptation: none was needed and none was made; every reference inside the delivered unit resolves inside it, so no unresolved reference or citation is produced. Printed slips kept literal: no printed word, symbol, hypothesis or number of the counted statement or of its proof was altered. Layout and class: the article's own class is \texttt{amsart} with packages including amssymb, amsmath, amsthm, amsfonts, mathtools, bbm, mathrsfs, enumitem, graphics, float, multirow, ulem, hyperref, comment; the wrapper is the lane's pinned \texttt{amsart} editorial wrapper at 10pt on A4, which restates the theorem-like environments the retained text needs and adds the article's own macro definitions that the retained text needs (\texttt{\textbackslash geq}, \texttt{\textbackslash leq}), with extra wrapper packages (bm, bbm, dsfont, xspace, cancel, mathtools). The delivered numbering is the unit's own. Assets: no retained span contains a figure environment, a graphics-inclusion command, a PDF-page inclusion, a table or a TikZ picture; no image file is copied into this package, the package declares no asset dependency and no image file is redistributed with it. Licence: version arXiv:2503.00990v2 of this article is distributed under the Creative Commons Attribution 4.0 International licence, as recorded in the arXiv licence metadata for this exact version and shown on the abstract page https://arxiv.org/abs/2503.00990v2. A full rights notice is delivered at licenses/arxiv-2503.00990-CC-BY-4.0.txt. Characters: every retained excerpt is pure ASCII, so the delivered engine, XeLaTeX with its default Latin Modern text font, reports no missing character.
\begin{lemma} \label{lemma:nested sequence}
Let $A \subset Q_{n,q}$ be such that $\langle A \rangle =Q_{n,q}$. Then there is a nested sequence $Q_0=Q_{i_1}^{x_{i_1}} \subset Q_{i_2}^{x_{i_2}} \subset \cdots \subset Q_{i_t}^{x_{i_t}}=Q_{n,q}$, of internally spanned subcubes with respect to $A$, where $2i_j+2 \geq i_{j+1}$ for all $j$, $ 0 \leq j \leq t-1$. Furthermore, for $j \geq 2$ each subcube $Q_{i_j}^{x_{i_j}}$ is spanned by two internally spanned cubes, namely by $Q_{i_{j-1}}^{x_{i_{j-1}}}$ and a subcube $Q_{m_{j-1}}$ with $m_{j-1} \leq i_{j-1}$ which is not a member of the sequence. 
\end{lemma}

\begin{lemma} \label{lemma:ST3}
Let $k,l \in \mathbb{N}_0$, $n=k+l+2$. Let $S=[*]^k[0]^{l+2}$ and $T=[0]^k[*]^lab$. With $A_0=S \cup T,$ where $a,b \neq 0$, then we have
$$A_1 \supset \{[0]^{k+l}0b,[0]^{k+l}a0,[0]^{k+l}ab,[0]^{k+l}00\},$$
$$A_2 \supset \{[0]^{k+l}0*,[0]^{k+l}a*,[0]^{k+l}*0,[0]^{k+l}*b\},$$
$$A_3 \supset \{[0]^{k+l}**\},$$
and for all $t \geq 2$,

\begin{align}
 A_t & \supset \{[0]^{k+l}**\}\label{eq:0b}
 \bigcup\{x \in \{0,1,\dots,q-1\}^n:  1 \leq \left \lVert x \right \rVert^{0b} \leq t-1\}\\ 
 \label{eq:a0}
&\bigcup \{x \in \{0,1,\dots,q-1\}^n:  1 \leq \left \lVert x \right \rVert^{a0} \leq t-1\}\\ 
\label{eq:0i}
&\bigcup_{i \neq b} \{x \in \{0,1,\dots,q-1\}^n:  1 \leq \left \lVert x \right \rVert^{0i} \leq t-2\} \\  
\label{eq:jo}
&\bigcup_{j \neq a} \{x \in \{0,1,\dots,q-1\}^n:  1 \leq \left \lVert x \right \rVert^{j0} \leq t-2\} \\ 
\label{eq:ai}
&\bigcup_{i \neq 0} \{x \in \{0,1,\dots,q-1\}^n:  1 \leq \left \lVert x \right \rVert^{ai} \leq t-2\} \\ 
\label{eq:jb}
&\bigcup_{j \neq 0} \{x \in \{0,1,\dots,q-1\}^n:  1 \leq \left \lVert x \right \rVert^{jb} \leq t-2\} \\ 
\label{eq:cd}
&\bigcup_{\substack{c \neq 0,a \\ d \neq 0,b}} \{x \in \{0,1,\dots,q-1\}^n:  1 \leq \left \lVert x \right \rVert^{cd} \leq t-3\}. 
 \end{align}




\end{lemma}

\begin{lemma} \label{lemma:ST6}
Let $k,l \in \mathbb{N}_0$, $n=k+l+2$. Let $S=[*]^k[0]^{l+2}$ and $T=[0]^k[*]^lab$ where $a,b \neq 0$. With $A_0=S \cup T,$ then we have the following. 

For $t \geq 1$,
$$A_t \cap \{x \in \{0,1,\dots,q-1\}^n: \left \lVert x \right \rVert_k^{0b} \geq t\}=\emptyset,$$
and 
$$A_t \cap \{x \in \{0,1,\dots,q-1\}^n: \left \lVert x \right \rVert_k^{a0} \geq t\}=\emptyset.$$

For $j \notin \{b,0\}$ and $t \geq 2$,
$$A_t \cap \{x \in \{0,1,\dots,q-1\}^n: \left \lVert x \right \rVert_k^{0j} \geq t-1\}=\emptyset.$$

For $i \notin \{a,0\}$ and $t \geq 2$,
$$A_t \cap \{x \in \{0,1,\dots,q-1\}^n: \left \lVert x \right \rVert_k^{i0} \geq t-1\}=\emptyset.$$

For $j \neq 0$ and $t \geq 2$,
$$A_t \cap \{x \in \{0,1,\dots,q-1\}^n: \left \lVert x \right \rVert_k^{aj} \geq t-1\}=\emptyset.$$

For $i \neq 0$ and $t \geq 2$,
$$A_t \cap \{x \in \{0,1,\dots,q-1\}^n: \left \lVert x \right \rVert_k^{ib} \geq t-1\}=\emptyset.$$

For $i \notin \{a,0\}$, $j \notin \{b,0\}$, and $t \geq 3$, 
$$A_t \cap \{x \in \{0,1,\dots,q-1\}^n: \left \lVert x \right \rVert_k^{ij} \geq t-2\}=\emptyset.$$

\end{lemma}

\begin{lemma} \label{lemma:summary}
For $x,y \in \{0,1,\cdots,q-1\}^n$ such that dim($x$)=$k$, dim($y$)=l, dim($x \vee y$)=$m$, where $k,l <m \leq n$, $d(x,y)=d \leq 2$, and $|\{i: x_i=y_i=* \}|=p$, then the infection time of $Q^x \cup Q^y$ in $Q_{n,q}$ is:

For $d=0$, $$T_{Q_{n,q}}(Q^x \cup Q^y)=m-p-1.$$

For $d=1$, $$T_{Q_{n,q}}(Q^x \cup Q^y)=m-p.$$

For $d=2$, $$T_{Q_{n,q}}(Q^x \cup Q^y)=m-p+1.$$
\end{lemma}

\begin{theorem}
For $q \geq 3$, $M_q(0)=0$ $M_q(1)=1$, and $M_q(2)=3$. Moreover,  for all $n \geq 3$
$$M_3(n)-M_3(n-3)=2n-1,$$
and for $q \geq 4$
$$M_q(n)-M_q(n-3)=2n.$$
\end{theorem}

\begin{proof}
It is obvious that $M_{q}(0)=0$ and $M_q(1)=1$. Now consider the case where $n=2$. W.l.o.g we can assume $00 \in A_0$ and first consider $A_0=\{00,ab\}$ where $a, b\neq 0$. It takes 3 steps to infect every vertex in $Q_q(2)$. Another initially infected set that we need to consider is $A_{0}=\{00,a0,0b\}$ where $a,b \neq 0$ and again it takes 3 steps to percolate. These are the only sets of initially infected vertices that we need to consider since adding any other vertex to these sets will not increase the percolation time and the process can not be initiated if any vertex is removed from these sets. 

We will first prove the lower bound of $M_{q}(n)$. 

Case 1. Let $A^{n-2}$ be a set that internally spans $Q^x_{n-2}$ for $x=[\ast]^{n-2}00$ in time $M_q(n-2)$ such that $[0]^n$ is infected at time $M_q(n-2)$. Now let $A'=A^{n-2} \cup [0]^{n-2}11$ and assume that we initially infect every vertex in $A'$. Note that $\langle A' \rangle=Q_{n,q}$ and before the time $M_q(n-2)$, none of the vertex in $[0]^{n-2}\ast \ast$ has interaction with the infection process in $[\ast]^{n-2}00$. 

Now assume that every vertex in $[\ast]^{n-2}00$ has been infected. Then it takes exactly $n+1$ additional steps to infect the rest of the vertices in $Q_{n,q}$. Indeed, by Lemma \ref{lemma:summary}, we have $d([\ast]^{n-2}00, [0]^{n-2}11)=2$, $|\{i:x_i=y_i=\ast\}|=0$ where $x=[\ast]^{n-2}00$ and $y=[0]^{n-2}11$, and $\langle [\ast]^{n-2}00, [0]^{n-2}11)=2\rangle= Q_{n,q}$. Therefore, 
$$T_{Q_{n,q}}(A')=M_q(n-2)+n+1.$$

Case 2. Let $A^{n-3}$ be a set that internally spans $Q_{n-3}^x$ for $x=[\ast]^{n-3}000$ in time $M_q(n-3)$ such that $[0]^n$ is infected at time $M_q(n-3)$. Let $A'=[\ast]^{n-3}000 \cup [0]^{n-3}110 \cup [2]^n$ and assume we initially infect every vertex in $A'$. Note that 
$\langle A' \rangle=Q_{n,q}$ and the set of sites infected after $M_q(n-3)$ steps is $[\ast]^{n-3}000 \cup [0]^{n-3}110 \cup [2]^n$.

After $M_q(n-3)+n-1$ steps, we claim that $[2]^n$ is the only vertex infected in $\cup_{j \neq 0} ([\ast]^{n-1}j)$. Indeed, the vertices $y$ in $[\ast]^{n-1}0$ that are at distance 2 from $[2]^n$ satisfy for $i \neq 2$ 
\begin{align} \label{eq:j20}
\left \lVert y \right \rVert^{i20}= n-3,  
\end{align} 

for $j \neq 2$
\begin{align}\label{eq:2j0}
    \left \lVert y \right \rVert^{2j0}= n-3,
\end{align}

 \begin{align} \label{eq:220}
\left \lVert y \right \rVert^{220}= n-4 
 \end{align}

or 
\begin{align} \label{eq:220'}
    \left \lVert y \right \rVert^{220}=n-3.
\end{align}
 

By Lemma \ref{lemma:ST3} and Lemma \ref{lemma:ST6}, none of those vertices $y$ satisfying (\ref{eq:j20})--(\ref{eq:220'}) is infected at time $ \leq M_q(n-3)+n-2$. In fact some of those vertices $y$ are infected exactly at time $M_q(n-3)+n-1$ and the rest of those vertices will be infected at time $M_q(n-3)+n$. Therefore, at time $M_q(n-3)+n-1$, $[2]^n$ is the only vertex infected in $\cup_{j \neq 0} ([\ast]^{n-1}j)$. Now we assume that every vertex in $[\ast]^{n-1}0 \cup [2]^n$ has been infected, then by Lemma \ref{lemma:summary}, it takes exactly $n$ steps to infected the entire $Q_{n,q}$ since $d([\ast]^{n-1}0,[2]^n)=1$, $|\{i:x_i =y_i=\ast\}|=0$ where $x=[\ast]^{n-1}0$ and $y=[2]^n$, and $\langle[\ast]^{n-1}0 \cup [2]^n \rangle=Q_{n,q}$.

Therefore, we have 
$$T_{Q_{n,q}}(A') \geq M_q(n-3)+2n-1.$$

Additionally we would like to show that for $q \geq 4$,
\begin{align}\label{eq:T(Qn)}
    T_{Q_{n,q}}(A') \geq M_q(n-3)+2n.
\end{align} 

In order to prove this tighter bound we need to further analyze the process right after $M_q(n-3)$ steps.

Now we need to show the following claim for $q \geq 4$. 
\begin{claim} \label{claim:subset I}
    For $ 1 \leq k \leq n-1$, every $k$ dimensional cube in $[\ast]^{n-1}2$ will be fully infected at time $M_q(n-3)+n+k$ and at time $\leq M_q(n-3)+n+k-1$ none of the $k$ dimensional cubes in $[\ast]^{n-1}2$ has been fully infected. In particular, the vertices $x$ with $x_n=2$ satisfying the following condition, denoted as Condition($k$),are not infected at time  $\leq M_q(n-3)+n+k-1$:
    
     there exists a subset $I \subset [n-1]$ of size at least $k$ such that 
\begin{align*}
x_i &=\begin{cases} 
  \notin \{0,2\}& \text{if $i \in I$ and $i \notin \{n-2,n-1\}$} \\
   \notin \{0,1,2\} & \text{if $i \in I$ and $ i \in \{n-2,n-1\}$} \\
   \in \{0,2\}& \text{if $i \notin I$ and $i \notin \{n-2,n-1\}$}\\
   \in \{0,1,2\} & \text{if $i \notin I$ and $i \in \{n-2,n-1\}$}
   \end{cases}
\end{align*}



    
    For $l \neq 0,2$ and  for $ 1 \leq k \leq n-1$, none of the $k$ dimensional cubes in $[\ast]^{n-1}l$ has been fully infected at time $ \leq M_q(n-3)+n+k$. In particular, the vertices $x$ with $x_n=l$ satisfying Condition($k$) are not infected at time $ \leq M_q(n-3)+n+k$.
    
\end{claim} 
\begin{proof}
Let $k=1$. Note that by Lemma \ref{lemma:ST3} and Lemma \ref{lemma:ST6}, after $M_q(n-3)+n-1$ steps every vertex $x$ in $[\ast]^{n-1}0$ has been infected except for the vertices $x$ satisfying 
$$\left \lVert x \right \rVert^{ij0}= n-3$$
where $i,j \notin \{0,1\}$
and after one additional step, the cube $[\ast]^{n-1}0$ has been fully infected. 

Again $[2]^n$ is the only vertex infected in $\cup_{j \neq 0} ([\ast]^{n-1}j)$ at time $M_q(n-3)+n-1$ and after one more step there are no infected vertices in $\cup_{j \notin \{0,2\}} ([\ast]^{n-1}j)$ besides $[2]^n$. The following vertices are infected in $[\ast]^{n-1}2$ at time $M_{q}(n-3)+n$:
  \begin{gather*}
 [2]^{n-3}122 \quad [2]^{n-3}022\\
[2]^{n-3}202 \quad [2]^{n-3}212
   \end{gather*}
and 
\begin{gather*}
0[2]^{n-1}\\
202^{n-2}\\
2202^{n-3}\\
...\\
[2]^{n-4}0[2]^3.
\end{gather*}

Indeed the vertex $[2]^{n-3}122$ has two neighbors infected at time $ \leq M_q(n-3)+n-1$, namely $[2]^n$ and $[2]^{n-3}120$ and it is similar for other vertices listed above. 


At time $M_q(n-3)+n+1$, every $1$ dimensional cube in $[\ast]^{n-1}2$ is fully infected since the vertices in a $1$ dimensional cube are either infected at time $M_q(n-3)+n$ or having at least $2$ neighbors infected at time $M_q(n-3)+n+1$. 

Now consider the vertices in $\cup_{j \notin \{0,2\}} ([\ast]^{n-1}j)$. At time $M_q(n-3)+n+1$, the following vertices are infected:
$$\cup_{j \notin \{0,2\}}[2]^{n-1}j$$
since the vertex $[2]^{n-1}j$ has two neighbors infected at time $\leq M(n-3)+n$ namely $[2]^n$ and $[2]^{n-1}0$. Additionally, at time $M_q(n-3)+n+1$, for all $l \notin \{0,2\}$
  \begin{gather*}
 [2]^{n-3}12l \quad [2]^{n-3}02l\\
[2]^{n-3}20l \quad [2]^{n-3}21l
   \end{gather*}
and 
\begin{gather*}
0[2]^{n-2}l\\
202^{n-3}l\\
2202^{n-4}l\\
...\\
[2]^{n-4}0[2]^2l
\end{gather*}
are infected. It is easy to see that the rest of the vertices in $\cup_{j \notin \{0,2\}} ([\ast]^{n-1}j)$ are not infected at time $M_{q}(n-3)+n+1$.

Therefore, the statement is true for $k=1$.

Now assume the statement is true for $k \geq 1$. Consider a vertex $x$ with $x_n=2$ satisfying Condition($k+1$). By the induction hypothesis this vertex $x$ has only one neighbors at time $M_q(n-3)+n+k-1$, namely, $x_1x_2\cdots x_{n-1}0$. 

Indeed, let us choose a particular $i \in [n-1]$ where $i \notin \{n-1,n-2\}$ and $i \notin I$ and w.l.o.g we can assume that $x_i=0$. Consider a neighbor $y$ of $x$ where $y_j=x_j$ for $j \neq i$ and $y_i=2$. By the induction hypothesis, the vertex $y$ is not infected at time $M_q(n-3)+n+k-1$. Consider another neighbor $z$ of $x$ where $z_j=x_j$ for $j \neq i$ and $z_i=a$ where $a \notin \{0,2\}$. By the induction hypothesis, $z$ is not infected at time $M_q(n-3)+n-k-1$ since the vertex $z$ satisfies Condition($k+2$). The analysis is very similar for $i \in \{n-1,n-2\}$ and $i \notin I$.  

Now we choose a particular $i \in [n-1]$ where $i \notin \{n-1,n-2\}$ and $i \in I$. Consider a neighbor $y$ of $x$ where $y_j=x_j$ for $j \neq i$ and $y_i \notin \{x_i,0,2\}$. By the induction hypothesis $y$ is not infected at time $M_q(n-3)+n+k-1$ since the vertex $y$ satisfies Condition($k+1$). Consider another neighbor $z$ of $x$ with $z_j=x_j$ for $j \neq i$ and $z_i=a$ where $a \in \{0,2\}$. By the induction hypothesis the vertex $z$ is not infected at time $\leq M_q(n-3)+n+k-1$ since $z$ satisfies Condition($k$). The analysis is very similar for $i \in \{n-1,n-2\}$ and $i \notin I$.

Now consider a neighbor $y$ of $x$ where $y_j=x_j$ for $j \neq n$ and $y_n \notin \{0,2\}$.By the induction hypothesis, the vertex $y$ is not infected at time $M_q(n-3)+n-k$ since $y$ satisfies Condition($k$). 

Consider a vertex $x$ with $x_n=l$ satisfying Condition($k+1$). By the induction hypothesis this vertex $x$ has only one infected neighbor at time $M_q(n-3)+n+k$, namely, $x_1x_2\cdots x_{n-1}0$. The analysis is very similar to the case where $x_n=2$. 
\end{proof}
It is easy to see that claim \ref{claim:subset I} suffices to show (\ref{eq:T(Qn)}).

\vspace{5mm}
Now let us try to prove the upper bound of $M_{q}(n)$. Let $A$ be a set of initially infected vertices spanning the hypercube $Q_{n,q}$ by the process. 

Let  $Q_0=Q_{i_1}^{x_{i_1}} \subset Q_{i_2}^{x_{i_2}} \subset \cdots \subset Q_{i_t}^{x_{i_t}}=Q_{n,q}$ be a nested sequence of internally spanned subcubes with respect to $A$ and $Q_{m_1}^{Z_{m_1}}$,$Q_{m_2}^{Z_{m_2}}$,$\cdots$,$Q_{m_{t-2}}^{Z_{m_{t-2}}}$,$Q_{m_{t-1}}^{Z_{m_{t-1}}}$ be the cubes that marge with cubes $Q_{i_j}^{x_{i_j}}$ as described in Lemma \ref{lemma:nested sequence}. Assume w.l.o.g both sequences are of maximal lengths. 

W.l.o.g, assume that $A$ is minimal (under containment) set spanning $Q_{n,q}$. Note that $i_{t-1} \leq n$. We will proceed by analyzing different cases based on the the values of $i_{t-1}$ and $i_{t-2}$. By Lemma \ref{lemma:nested sequence}, for all $ 1 \leq j \leq t-1$ we have $i_j \geq m_j$

Case 1. If $i_{t-1} \leq n-2$, then after at most $M_q(i_{t-1}) \leq M_q(n-2)$ steps, both $Q_{i_{t-1}}^{x_{i_{t-1}}}$ and $Q_{m_{t-1}}^{z_{m_{t-1}}}$ are fully infected. Since $\langle Q_{i_{t-1}}^{x_{i_{t-1}}} \cup Q_{m_{t-1}}^{z_{m_{t-1}}} \rangle=Q_{n,q}$, after at most additional $n+1$ steps every vertex in $Q_{n,q}$ will be infected by Lemma \ref{lemma:summary}. Indeed, we have dim($Q_{i_{t-1}}^{x_{i_{t-1}}} \vee Q_{m_{t-1}}^{z_{m_{t-1}}}$)=n and consider the worst case where $d(x,y)=2$, and $|\{i:x_i=y_i=\ast \}|=0$ with $x=Q_{i_{t-1}}^{x_{i_{t-1}}}$ and $y=Q_{m_{t-1}}^{z_{m_{t-1}}}$.
Therefore, we have 
$$T_{Q_{n,q}}(A) \leq M(n-2)+n+1$$

Case 2. If $i_{t-1}=n-1$ and $i_{t-2}=n-2$, then there is a vertex $v \in A\cap Q_{m_{t-2}}^{z_{m_{t-2}}}$ such that $d(x_{t-2},v)=1$ and $\langle Q_{i_{t-2}}^{x_{i_{t-2}}} \cup v \rangle=Q_{i_{t-1}}^{x_{i_{t-1}}}.$ Moreover, there exists a vertex $w \in A \cap Q_{m_{t-1}}^{z_{m_{t-1}}}$ such that $\langle Q_{i_{t-1}}^{x_{i_{t-1}}} \cup w \rangle=Q_{n,q}.$ 

Note that $d(x_{i_{t-2}},w)=1$ or $d(x_{i_{t-2}},w)=2$. If $(x_{i_{t-2}},w)=2$, then $\langle Q_{i_{t-2}}^{x_{i_{t-2}}} \cup w \rangle=Q_{n,q},$ which contradicts the minimality of $A$ since $\langle A \backslash \{v\} \rangle =Q_{n,q}$. Thus $d(x_{i_{t-2}},w)=1$.

W.l.o.g we can assume that
$$x_{i_{t-2}}=[\ast]^{n-2}00,$$
$$x_{i_{t-2}} \vee v=[\ast]^{n-1}0,$$
and 
$$x_{i_{t-2}} \vee w=[\ast]^{n-2}0\ast.$$

It is easy to see that after at most $M_q(n-2)$ steps the cube $Q_{i_{t-2}}^{x_{i_{t-2}}}$ is fully infected. After at most $(n-1)$ additional steps both $Q_{n-1}^{x_{i_{t-2}} \vee v}$ and $Q_{n-1}^{x_{i_{t-2}} \vee w }$ are fully infected by Lemma \ref{lemma:summary}. After 2 more steps every vertex in $Q_{n,q}$ will be infected again by Lemma \ref{lemma:summary}. Therefore, we have
$$T_{Q_{n,q}}(A) \leq M_q(n-2)+n-1+2 =M_q(n-2)+n+1.$$

Case 3. If $i_{t-1}=n-1$, $i_{t-2} \leq n-3$, and $d({x_{i_{t-2}}},z_{m_{t-2}}) \leq 1$, then after $M(n-3)$ steps both $Q_{i_{t-2}}^{x_{i_{t-2}}}$ and $Q_{m_{t-2}}^{z_{m_{t-2}}}$ are fully infected. Since $\langle Q_{i_{t-2}}^{x_{i_{t-2}}} \cup Q_{m_{t-2}}^{z_{m_{t-2}}} \rangle=Q_{i_{t-1}}^{x_{i_{t-1}}}$, after at most $i_{t-1}$ more steps the cube $Q_{i_{t-1}}^{x_{i_{t-1}}}$ is fully infected by Lemma \ref{lemma:summary}. 

Since $i_{t-1}=n-1$, $d(x_{i_{t-1}},Z_{m_{t-1}}) \leq 1$. Again by Lemma \ref{lemma:summary} after at most $n$ more steps every vertex in $Q_{n,q}$ will be infected. Therefore, we have 
$$T_{Q_{n,q}}(A) \leq M_q(n-3)+n-1+n=M_q(n-3)+2n-1.$$

Case 4. If $i_{t-1}=n-1$,$i_{t-2} \leq n-3$, $d(x_{i_{t-2}},z_{m_{t-2}})=2$, and $m_{t-2}=n-3$, then after at most $M(n-3)$ steps both $Q_{i_{t-2}}^{x_{i_{t-2}}}$ and $Q_{m_{t-2}}^{z_{m_{t-2}}}$ are fully infected. Note that $\langle Q_{i_{t-2}}^{x_{i_{t-2}}} \cup Q_{m_{t-2}}^{z_{m_{t-2}}} \rangle=Q_{i_{t-1}}^{x_{i_{t-1}}}$ and since $i_{t-2} \geq m_{t-2}$ then $i_{t-2}=n-3$. After at most 3 additional steps the cube $Q_{i_{t-1}}^{x_{i_{t-1}}}$ is infected. Indeed observe that $|\{i: x_i=y_i=\ast\}| =n-3$ where $x= Q_{i_{t-2}}^{x_{i_{t-2}}}$ and $y= Q_{i_{t-2}}^{x_{i_{t-2}}}$ since one of the coordinates must be the same for $x$ and $y$. since $d(x,y)=2$ and $\text{dim}(x \vee y)=n-1$, $T_{Q_{n,q}}(Q_{i_{t-2}}^{x_{i_{t-2}}} \cup Q_{m_{t-2}}^{z_{m_{t-2}}})=3$ by Lemma \ref{lemma:summary}.

Since $i_{t-1}=n-1, d(x_{i_{t-1}},z_{m_{t-1}}) \leq 1$. By lemma \ref{lemma:summary} after at most $n$ more steps every vertex on $Q_{n,q}$ will be infected. Therefore, we have 
$$T_{Q_{n,q}}(A) \leq M_q(n-3)+3+n.$$

Additionally we would like to show 
$$T_{Q_{n,3}}(A) \leq M_3(n-3)+2+n.$$
W.l.o.g, we can assume that 
$$x_{i_{t-1}}=[\ast]^{n-1}0,$$
$$x_{i_{t-2}}=[\ast]^{n-3}000,$$
and 
$$z_{m_{t-2}}=[\ast]^{n-3}ij0,$$
where $i \neq 0$ and $j \neq 0$.

Again after $M_3(n-3)$ steps, both $Q_{i_{t-2}}^{x_{i_{t-2}}}$ and $Q_{m_{t-2}}^{z_{m_{t-2}}}$ are fully infected. After at most 2 additional steps every vertex in $Q_{i_{t-1}}^{x_{i_{t-1}}}$ satisfying $[\ast]^{n-3}kl0$ for all $k,l\in \{0,1,2\}$ but $k \notin \{0,i\}$ and $l \notin \{0,j\}$. 

There is a vertex $y \in A$ such that $A \in [\ast]^{n-1}l$ where $l \neq 0$. Otherwise the vertices not in the cube $Q_{i_{t-1}}^{x_{i_{t-1}}}$ will not be infected. Let $y=y_1y_2 \cdots y_{n-1}l$.

Therefore, after $M_3(n-3)+2$ steps, the vertices $v$ and $w$ in $Q_{i_{t-1}}^{x_{i_{t-1}}}$ are infected, where 
$$v=y_1y_2\cdots y_{n-3}\;i\;y_{n-1}0$$
and 
$$w=y_1y_2\cdots y_{n-3}\;0\; y_{n-1}0.$$
Assume $y_{n-2} \notin \{0,i\}$. Then after $M_3(n-3)+n$, vertices $v^*$ and $w^*$ in $[\ast]^{n-1}l$ are infected, where 
$$v^*=y_iy_2 \cdots y_{n-3}\;i\;y_{n-1}l$$
and 
$$w^*=y_iy_2 \cdots y_{n-3}\;0\;y_{n-1}l. $$
Therefore, after $M_3(n-3)+3$, $[\ast]^{n-1}0$ and $y_1y_2\cdots y_{n-3}\ast y_{n-1}l$ have been infected. By Lemma \ref{lemma:summary} after at most $n-1$ steps every vertex in $Q_{n,3}$ will be infected. 

Now assume $y_{n-2}=0$. If $y_{n-1} \notin \{0,j\}$, similarly it can be shown that after $M_3(n-3)+n+2$ every vertex in $Q_{n,3}$ will be infected. Thus we can assume $y_{n-1} \in \{0,j\}$. First let $y_{n-1}=0$. Note that after $M_3(n-3)+n-1$ steps, the vertices $u$ and $z$ in $Q_{i_{t-1}}^{x_{i_{t-1}}}$ are infected where 
$$u=y_1y_2 \cdots y_{n-3}k00$$
with $k \notin\{i,0\}$
and 
$$z=y_1y_2 \cdots y_{n-3}i00.$$
Then after one more step, the vertices $u^*$ and $z^*$ in $[\ast]^{n-1}l$ are infected, where 
$$u^*=y_1y_2 \cdots y_{n-3}k0l$$
and 
$$z^*=y_1y_2 \cdots y_{n-3}i0l.$$
Therefore, after $M_3(n-3)+3$ steps, $[\ast]^{n-1}0$ and $y_1y_2\cdots y_{n-3}\ast 0l$ have been infected. By Lemma \ref{lemma:summary} after at most $n-1$ steps every vertex in $Q_{n,3}$ will be infected. 

Due to symmetry for other values of $y_{n-2}$ and $y_{n-1}$ the argument is the same. 

Therefore, we have 
$$T_{Q_{n,3}} \leq M_3(n-3)+n+2.$$


Case 5. If $i_{t-1}=n-1$,$i_{t-2} \leq n-3$, $d(x_{i_{t-2}},z_{m_{t-2}})=2$, and $m_{t-2} < n-3$, after at most $M_q(n-3)$ steps, both $Q_{i_{t-2}}^{x_{i_{t-2}}}$ and $Q_{m_{t-2}}^{z_{m_{t-2}}}$ are fully infected. After at most $n$ more steps the cube $Q_{i_{t-1}}^{x_{i_{t-1}}}$ is fully infected by Lemma \ref{lemma:summary}. Then after at most $n$ more 
steps every vertex on $Q_{n,q}$ is infected again by Lemma ~\ref{lemma:summary} since $d(x_{i_{t-2}},z_{m_{t-2}}) \leq 1$ and $\langle Q_{i_{t-1}}^{x_{i_{t-1}}} \cup Q_{m_{t-1}}^{z_{m_{t-1}}} \rangle=Q_{n,q}.$ Therefore, we have for $q \geq 3$
$$T_{Q_{n,q}}(A) \leq M_q(n-3)+n+n=M_q(n-3)+2n.$$

Additionally we would like to show that 
$$T_{Q_{n,3}}(A) \leq M_3(n-3)+2n-1.$$

In order to show this tighter bound we need to analyze the process after $M_3(n-3)$ steps in a more careful way. Assume that $q=3$. 

W.l.o.g we can assume that 
$$x_{i_{t-1}}=[\ast]^{n-1}0,$$
$$x_{i_{t-2}}=[\ast]^{i_{t-2}}[0]^{n-i_{t-2}},$$
and 
$$z_{m_{t-2}}=[\ast]^p[0]^{i_{t-2}-p}[\ast]^{m_{t-2}-p}ij0,$$
where $i \neq 0$ and $j \neq 0$.

Again after at most $M(n-3)$ steps both $Q_{i_{t-2}}^{x_{i_{t-2}}}$ and $Q_{m_{t-2}}^{z_{m_{t-2}}}$ are fully infected. After at most $n-1$ additional steps, by Lemma \ref{lemma:ST3} and Lemma \ref{lemma:ST6}, vertices $x$ in $Q_{i_{t-1}}^{x_{i_{t-1}}}$ satisfying $ 1 \leq \left \lVert x \right \rVert^{kl0} \leq n-3$  for all $k,l\in \{0,1,2\}$ but $k \notin \{0,i\}$ and $l \notin \{0,j\}$ are infected.

There is a vertex $y \in A$ such that $A \in [\ast]^{n-1}l$ where $l \neq 0$. Otherwise the vertices not in the cube $Q_{i_{t-1}}^{x_{i_{t-1}}}$ will not be infected. Let $y=y_1y_2 \cdots y_{n-1}l$. 

Assume that $\sum_{i=1}^{n-3}\mathbb{I}_{\{y_i \neq 0\}} \geq 1$. Therefore, after $M_3(n-3)+n-1$ steps, the vertices $v$ and $w$ in $Q_{i_{t-1}}^{x_{i_{t-1}}}$ are infected, where 
$$v=y_1y_2\cdots y_{n-3}\;i\;y_{n-1}0$$
and 
$$w=y_1y_2\cdots y_{n-3}\;0\; y_{n-1}0.$$
Assume $y_{n-2} \notin \{0,i\}$. Then after $M_3(n-3)+n$, vertices $v^*$ and $w^*$ in $[\ast]^{n-1}l$ are infected, where 
$$v^*=y_iy_2 \cdots y_{n-3}\;i\;y_{n-1}l$$
and 
$$w^*=y_iy_2 \cdots y_{n-3}\;0\;y_{n-1}l. $$
Therefore, after $M(n-3)+n$, $[\ast]^{n-1}0$ and $y_1y_2\cdots y_{n-3}\ast y_{n-1}l$ have been infected. By Lemma \ref{lemma:summary} after at most $n-1$ steps every vertex in $Q_{n,3}$ will be infected. 

Now assume $y_{n-2}=0$. If $y_{n-1} \notin \{0,j\}$, similarly it can be shown that after $M_3(n-3)+2n-1$ every vertex in $Q_{n,3}$ will be infected. Thus we can assume $y_{n-1} \in \{0,j\}$. First let $y_{n-1}=0$. Note that after $M_3(n-3)+n-1$ steps, the vertices $u$ and $z$ in $Q_{i_{t-1}}^{x_{i_{t-1}}}$ are infected where 
$$u=y_1y_2 \cdots y_{n-3}k00$$
with $k \notin\{i,0\}$
and 
$$z=y_1y_2 \cdots y_{n-3}i00.$$
Then after one more step, the vertices $u^*$ and $z^*$ in $[\ast]^{n-1}l$ are infected, where 
$$u^*=y_1y_2 \cdots y_{n-3}k0l$$
and 
$$z^*=y_1y_2 \cdots y_{n-3}i0l.$$
Therefore, after $M_3(n-3)+n$ steps, $[\ast]^{n-1}0$ and $y_1y_2\cdots y_{n-3}\ast 0l$ have been infected. By Lemma \ref{lemma:summary} after at most $n-1$ steps every vertex in $Q_{n,3}$ will be infected. 

Due to symmetry for other values of $y_{n-2}$ and $y_{n-1}$ the argument is the same. 

Therefore, we have 
$$T_{Q_{n,3}} \leq M_3(n-3)+2n-1.$$

Assume $y=[0]^{n-3}y_{n-2}y_{n-1}l$. Therefore, after $M_3(n-3)+n-1$ steps, as long as $n-1 \geq 2$, the vertices $v$ and $w$ in $Q_{i_{t-1}}^{x_{i_{t-1}}}$ are infected, where 
$$v=[0]^{n-3}\;i\;y_{n-1}0$$
and 
$$w=[0]^{n-3}\;0\; y_{n-1}0.$$
Indeed, the vertex $[0]^n$ and $[0]^{n-3}ij0$ are infected at time $\leq M_3(n-3)$. From now on the proof is identical to the case where $\sum_{i=1}^{n-3}\mathbb{I}_{\{y_i \neq 0\}} \geq 1$. 






%By Lemma \ref{lemma:ST3} after at most $M(n-3)+3$ steps the cube $[0]^{n-3}\ast \ast 0$ in %$Q_{i_{t-1}}^{x_{i_{t-1}}}$ is fully infected. After one more step the vertices  
%$$u=[0]^{n-3}ay_{n-1}l$$
%and 
%$$v=[0]^{n-3}by_{n-1}l$$
%are infected where $a \neq y_{n-2},b$ and $b \neq y_{n-2}$. Indeed the vertex $u$ has two infected neighbors, one of which is $y$ and the other is $[0]^{n-3}ay_{n-1}0$. Both of the neighbors are infected at time $\leq M(n-3)+3$. Therefore, after $M(n-3)+4$ steps, $[\ast]^{n-1}0$ and $y_1y_2\cdots y_{n-3}\ast 0l$ have been infected. By Lemma \ref{lemma:summary} after at most $n-1$ steps every vertex in $Q_{n,3}$ will be infected. 

%Therefore, we have 
%$$T_{Q_{n,3}} \leq M(n-3)+n+3.$$

Finally we have 
$$T_{Q_{n,3}} \leq M_3(n-3)+2n-1.$$

We need one more claim to conclude the proof. 
\begin{claim}
If $M_q(0)=0$,$M_q(1)=1$ and $M_q(2)=3$ and  
for $n \geq 3$ 
$$M_3(n)=\max\{M_3(n-2)+n+1,M_3(n-3)+2n-1\}$$
and for $q \geq 4$
$$M_q=\max\{M_q(n-2)+n+1,M_q(n-3)+2n\},$$ then 
$$M_3(n)=M_3(n-3)+2n-1$$
and for $q \geq 4$
$$M_q(n)=M_q(n-3)+2n.$$
\end{claim}
\begin{proof}
Assume that $q=3$. We have $M_3(3)=M_3(0)+5=5$ or $M_3(3)=M_3(1)+4=5$

$M_3(4)=M_3(1)+7=8$ or 
$M_3(4)=M_3(2)+5=8$.

$M_3(5)=M_3(2)+9=12$ or 
$M_3(5)=M_3(3)+6=11$.

Thus the lemma holds for $ 1 \leq n \leq 3$.

For $n \geq 3$, we assume that the lemma holds for $n,n+1,n+2$. We have 
\begin{align}
    M_3(n-3)& =\max\{M_3(n)+2(n+3)-1,M_3(n+1)+n+4\}\\
    &=\max\{M_3(n-3)+2n-1+2(n+3)-1, M_3(n-2)+2(n+1)-3+n+4\}\\
    &= \max{M_3(n-3)+4n+4,M_3(n-2)+3n+3}\\
   \label{eq:M(n-3+4n+4)} &=M_3(n-3)+4n+4\\ 
    &=M_3(n)-2n+1+4n+4\\
    &=M_3(n)+2n+5
\end{align}
where (\ref{eq:M(n-3+4n+4)}) follows from 
\begin{align*}
    M_3(n-3)+4n+4& =M_3(n-3)-2n+1+2n+5\\
    & \geq M_3(n-2)+n+1+2n+5\\
    & > M_3(n-2)+3n+3.
\end{align*}
For $q \geq 4$, the calculation is almost the same. 
\end{proof}
Therefore, we have our desired results.
\end{proof} 

\end{document}
